\section{Introduction}
\label{sec:introduction}

Retrieval can select passages about the right topic while missing a condition that determines which passage answers a query. Negation makes this failure particularly visible: two nearly identical passages can support opposite answers~\cite{weller-etal-2024-nevir}. After multiple retrievers contribute candidates, a reranker must both prefer the passage that satisfies the query and place it early enough to be used. These requirements motivate a joint examination of local preference and position in the returned list.

NevIR introduced contrastive passage pairs to diagnose negation sensitivity; its paired accuracy requires correct preferences for both associated queries~\cite{weller-etal-2024-nevir}. NevIR and its reproduction already show that improvements on negation-sensitive evaluation need not translate into improvements on broader retrieval benchmarks~\cite{weller-etal-2024-nevir,elsen_reproducing_2025}. We ask a more specific question within one evaluation environment: \emph{when the candidate set is fixed, do changes that improve the designated pairwise preference also improve the preferred passage's position among surrounding candidates?} A correct preference alone cannot answer this question: ranking the preferred passage 100th and its counterpart 101st satisfies the local ordering while leaving both far down the list.

We study this distinction in fixed pools formed by the union of dense, learned sparse, and BM25 retrieval results. All methods use the same retrieved passages, without additional retrieval or access to source text beyond those passages. We measure strict per-query preference for the officially designated preferred passage and its reciprocal rank in the full candidate list. Both metrics are conditional on pair coverage, which is reported separately. Background candidates lack relevance judgments, so these position metrics describe designated passages, not overall pool relevance or downstream answer quality.

On 371 covered confirmation queries, a feature-based ranker trained with pairwise supervision obtains 208 correct preferences versus 189 for fixed score fusion, yet has a lower mean reciprocal rank for the preferred passage (0.3169 versus 0.6017). A training variant with sampled background candidates and revised supervision restores this metric to 0.6133, but reduces correct preferences to 198. Thus, restoring position does not recover the stronger pairwise preference in this comparison. This joint intervention does not isolate the cause of recovery.

We also evaluate a bounded use of existing pointwise LLM reranking~\cite{zhuang-etal-2024-beyond}. A Qwen3-14B-AWQ configuration scores query--passage inputs independently within the fusion baseline's top 20. Fusion scores break ties between discrete judgments, while passages outside the window retain their order. The evaluation concerns the combined procedure: its scoring model, bounded window, and fusion prior all contribute to the resulting order.

Among the 371 covered confirmation queries, the first Qwen reranking run yields 280 correct preferences, compared with 189 for fusion and 257 for a BGE reranker using the same window. A separate fixed-configuration rerun yields 277; both runs are retained. These are observations on a previously used confirmation set, not independent Test evidence. The study contributes a within-pool diagnosis of the separation between pairwise preference and designated-passage position, together with an evaluation of training and bounded reranking interventions against both outcomes.
